\documentclass[12pt]{article}

% --- Geometry ---
\usepackage[margin=1.25in]{geometry}

% --- Font: Times New Roman family ---
\usepackage{mathptmx}
\usepackage[T1]{fontenc}
\usepackage[utf8]{inputenc}

% --- Typography ---
\usepackage{microtype}
\usepackage{setspace}
\onehalfspacing

% --- Math ---
\usepackage{amsmath}

% --- Tables ---
\usepackage{booktabs}
\usepackage{array}
\usepackage{tabularx}
\usepackage{makecell}
\usepackage{adjustbox}

% --- Figures ---
\usepackage{graphicx}
\graphicspath{{../../output/electoral_profile/boundary_permutation/}{../../output/diff/seed_001/plots/}}

% --- Captions ---
\usepackage[font=small,labelfont=bf,labelsep=period]{caption}

% --- References ---
\usepackage[authoryear,round]{natbib}

% --- Links ---
\usepackage[colorlinks=true,linkcolor=black,citecolor=black,urlcolor=blue,breaklinks=true]{hyperref}
\urlstyle{same}

% --- Section formatting ---
\usepackage{titlesec}
\titleformat{\section}{\normalfont\large\bfseries}{\thesection.}{0.5em}{}
\titleformat{\subsection}{\normalfont\normalsize\bfseries}{\thesubsection.}{0.5em}{}
\titlespacing*{\section}{0pt}{1.5ex plus 0.5ex minus 0.2ex}{0.8ex plus 0.2ex}
\titlespacing*{\subsection}{0pt}{1.2ex plus 0.3ex minus 0.2ex}{0.5ex plus 0.1ex}

% --- Abstract ---
\usepackage{abstract}
\renewcommand{\abstractnamefont}{\normalfont\bfseries}
\renewcommand{\abstracttextfont}{\normalfont\small}

% --- Footnote style ---
\renewcommand{\thefootnote}{\arabic{footnote}}

% --- Widow/orphan control ---
\widowpenalty=10000
\clubpenalty=10000

% ============================================================

\title{\textbf{Boundary Changes and Competitive Seat Exclusion in\\
Singapore's 2025 Redistricting: An Ensemble-Based Test}}

\author{Goh Cheng Arn David\thanks{Correspondence: \href{mailto:daveed@cs.toronto.edu}{daveed@cs.toronto.edu}.
Code and data: \url{https://github.com/davidcagoh/sg-redistricting-2025}.}}

\date{April 2026}

\begin{document}

\maketitle
\thispagestyle{empty}

\begin{abstract}
We test whether Singapore's 2025 electoral redistricting produced a statistically anomalous
allocation of competitively contested territory. Using a hypergeometric permutation test, we
examine the 114 subzones that changed constituency between the 2020 and 2025 general elections.
Among the 52 subzones originating from competitive 2020 constituencies (PAP vote share $<$55\%),
zero ended up in competitive 2025 constituencies---against an expected value of 3.2 under random
redistribution (Fisher's exact $p = 0.012$). The anomaly has a specific structure: the single new
competitive seat in 2025 (Jalan Kayu SMC) was seeded entirely from PAP stronghold population,
while the populations of the three dissolved competitive constituencies were absorbed into safe or
walkover districts. A complementary MCMC ensemble analysis finds that Singapore's constituency
boundaries are extreme outliers in planning-area cohesion and geometric compactness, but in a
direction consistent with administrative logic rather than partisan gerrymandering. The primary
mechanism producing Singapore's large seat-to-vote bonuses is the Group Representation
Constituency block-vote rule, not boundary drawing. The permutation test identifies a secondary
mechanism: the selective management of which communities become competitive through redistricting.
\end{abstract}

\vspace{1em}
\noindent\textit{Keywords:} redistricting, gerrymandering, Singapore, ensemble analysis, MCMC,
electoral geography, Group Representation Constituency

\clearpage

% ============================================================
\section{Introduction}

Singapore's 2025 general election produced a striking result: the ruling People's Action Party
(PAP) won 87 of 97 parliamentary seats---its highest proportion since 2015---while receiving 65.6\%
of the popular vote. The Workers' Party, the only credible opposition, held its two existing Group
Representation Constituencies (GRCs) but gained no new ground despite fielding a full slate and
running a disciplined campaign. The opposition seat count was the lowest in two decades of
contested elections.

The election took place less than three weeks after the Electoral Boundaries Review Committee
(EBRC) released its report. Singapore's boundaries are redrawn before every general election, and
2025 was no exception: 114 of 332 planning subzones changed their constituency allocation,
affecting approximately 1.3 million residents. Several high-profile marginal seats from
2020---constituencies where the PAP's margin was slim enough to be plausibly at risk---were
restructured or dissolved.

This paper asks a specific, narrow question: was the allocation of those 114 changed subzones
politically anomalous?

The question is narrow by design. We are not asking whether Singapore's electoral system is fair.
It is not, in the sense that it consistently produces large seat majorities from moderate vote
pluralities---but as we will show, that effect is structural, driven by the GRC block-vote
mechanism rather than by how boundaries are drawn. Nor are we asking whether boundaries were drawn
to create misshapen partisan districts in the manner of American partisan gerrymandering. The
evidence on that point runs the other way.

The specific question---were the changed subzones assigned in a way that is consistent with a
neutral redistribution?---can be answered with a statistical test. Using an approach inspired by
\citet{herschlag2020} ensemble analysis of North Carolina's congressional map, we generate a null
distribution for what competitive-origin subzone allocations would look like under random
redistribution, and we measure where Singapore's actual redistricting falls within that
distribution.

The answer is: at the 1.2nd percentile. Among 114 subzones that changed constituency, not one
from a competitive 2020 constituency ended up in a competitive 2025 constituency---when 3.2 would
be expected by chance. The exact probability of this result under the null hypothesis is $p =
0.012$.

The anomaly has a specific structure, which we describe in detail. The new competitive seat in
2025 was built from stronghold population, not from the populations of the three dissolved
competitive constituencies. Those populations were absorbed into safe or walkover districts. The
competitive electoral geography was, in effect, reset: the places that were battlegrounds in 2020
are no longer battlegrounds in 2025, and the new battlegrounds are in entirely different
communities.

What this finding cannot show is also worth stating up front. It cannot establish intent. The same
pattern is consistent with administrative population rebalancing that happened to relocate
competitive seats geographically. The test demonstrates that the pattern is statistically unlikely
under randomness; it cannot distinguish a deliberate strategy from coincidence. We return to this
distinction throughout.

% ============================================================
\section{Background}

\subsection{The Group Representation Constituency system}

Singapore's parliament uses a hybrid of single-member constituencies (SMCs) and Group
Representation Constituencies (GRCs). GRCs are multi-member constituencies where a team of 3 to
5 candidates stands together, and the team that wins a plurality takes all seats in the
constituency. Since 1988, GRCs have been justified primarily on ethnic representation grounds:
each GRC team must include at least one candidate from an ethnic minority group (Malay, Indian, or
Other), which the government argues prevents minority communities from being locked out of
parliament by simple-majority single-member voting.

The structural consequence of the GRC system is that it amplifies pluralities into seat bonuses.
A party that wins a 5-seat GRC by a single vote wins 5 seats; a party that loses by a single vote
wins zero. This winner-takes-all property, applied at the constituency level across seats that
range in size from 1 to 5, means that the seat-share-to-vote-share ratio is considerably higher
in Singapore than in a comparable single-member-plurality system.

In a standard single-member-plurality system, a 60\% national vote share would typically produce
approximately 60--70\% of seats, depending on the geographic concentration of the vote. In
Singapore, 60--65\% of the vote produces roughly 89\% of seats. This is not because the
boundaries are drawn to waste opposition votes in the way that American partisan gerrymandering
operates. It is because a 5-seat GRC won 55--45 produces a 5:0 seat split, not a 3:2 split.

\subsection{Singapore's seat--vote gap in numbers}

Table~\ref{tab:seatvote} documents the scale of seat magnification in recent elections.

\begin{table}[h]
\centering
\caption{Seat magnification in Singapore general elections.}
\label{tab:seatvote}
\begin{tabular}{lcccc}
\toprule
Year & PAP vote share & PAP seat share & Seat--vote gap & Efficiency gap \\
\midrule
2020 & 61.2\% & 89.3\% & $+$28.0 pp & 16.2\% \\
2025 & 65.6\% & 89.7\% & $+$24.1 pp & 7.3\% \\
\bottomrule
\end{tabular}
\end{table}

The efficiency gap---a measure developed to quantify partisan gerrymandering in the US
context \citep{mcghee2014}---was 16.2\% in 2020. For reference, the US Supreme Court in
\textit{Rucho v.\ Common Cause} (2019) examined maps with efficiency gaps in the 10--16\% range
as potential constitutional violations. Singapore's 2020 efficiency gap falls in that range,
though the institutional mechanism is different: it is the block-vote system, not boundary
manipulation, that drives it.

The votes-per-seat-won figure is more instructive. In 2020, the PAP won each seat with an average
of 16,928 votes; opposition parties won each seat with an average of 16,148 votes. The ratio is
0.95---essentially equal. This is not what partisan gerrymandering looks like. Gerrymandering, as
practised in the United States, works by over-concentrating opposition votes in a small number of
constituencies (packing) or spreading them thinly across many constituencies where they cannot win
(cracking). Both techniques would show up as a significant disparity in votes-per-seat. The
near-parity in Singapore means the seat--vote gap is produced elsewhere---in the multi-seat
block-vote mechanism, not in the geographic drawing of district lines.

\subsection{The ensemble approach}

\citet{herschlag2020} introduced the use of Markov Chain Monte Carlo (MCMC) sampling to evaluate
whether a specific redistricting plan is an outlier within the space of geometrically valid
alternatives. The key insight is that judging a single map as ``fair'' or ``gerrymandered''
requires a reference class: what would redistricting look like if it were done neutrally? MCMC
sampling generates that reference class by randomly exploring the space of plans satisfying stated
constraints---equal population, geographic contiguity, and whatever other neutral criteria
apply---and measuring where the actual plan falls in the resulting distribution.

Our implementation adapts this approach to Singapore. We run a 9,000-step ReCom MCMC chain on
Singapore's 332 subzones, treating them as the basic building blocks of electoral districts, with
equal-population and contiguity as the primary constraints. This generates an ensemble of 9,000
alternative ways to draw 33 constituency-sized units from the same subzones, producing a null
distribution against which the actual plans can be compared.

The boundary change permutation test (Section~\ref{sec:permutation}) follows the same logic but
asks a different question: not ``where does the actual boundary configuration fall among all
possible configurations?'' but rather ``given that 114 specific subzones changed constituency, was
the allocation of those subzones to destinations politically anomalous?'' This is a simpler test
with fewer modelling assumptions, which is why we treat it as the headline finding rather than the
MCMC results.

% ============================================================
\section{MCMC Ensemble Results}

Running the MCMC ensemble against both the 2020 and 2025 actual plans produces three striking
results and one uninformative one (Table~\ref{tab:mcmc}).

\begin{table}[h]
\centering
\caption{MCMC ensemble metrics. Ensemble drawn from 9,000 post-burn-in ReCom steps, $k=33$
districts, population tolerance 10\%.}
\label{tab:mcmc}
\begin{adjustbox}{max width=\textwidth}
\small
\begin{tabular}{lcccccc}
\toprule
Metric & \makecell{Actual\\2020} & \makecell{Actual\\2025} & \makecell{Ensemble\\mean} & \makecell{Ensemble\\range}
  & \makecell{2020\\pctile} & \makecell{2025\\pctile} \\
\midrule
Population deviation  & 1.223 & 1.223 & 0.070 & {[}0.069, 0.095{]} & 100th & 100th \\
\textbf{HDB towns split} & \textbf{12} & \textbf{12} & \textbf{13.0}
  & \textbf{{[}12, 13{]}} & \textbf{0th} & \textbf{0th} \\
\textbf{Planning area splits} & \textbf{30} & \textbf{28} & \textbf{34.3}
  & \textbf{{[}29, 41{]}} & \textbf{0th} & \textbf{0th} \\
\textbf{Compactness (PP)}    & \textbf{0.435} & \textbf{0.418} & \textbf{0.349}
  & \textbf{{[}0.296, 0.399{]}} & \textbf{100th} & \textbf{100th} \\
\bottomrule
\end{tabular}
\end{adjustbox}
\end{table}

The population deviation result is uninformative because the actual plan uses variable-size GRCs
(3--5 seats) while the ensemble models equal-population single-member districts. These are
structurally incompatible comparisons.

The three informative results are striking. \textit{HDB town cohesion}: the ensemble draws
overwhelmingly produce 13 HDB town splits (99.4\% of steps); the actual 2020 and 2025 plans
each split only 12 towns, tying the ensemble minimum. \textit{Planning area cohesion}: the 2020
actual plan achieves 30 planning area splits, below the 99th percentile of the ensemble (only
one of 9,000 randomly generated plans achieves fewer, at 29 splits); the 2025 plan achieves 28
splits, below the ensemble minimum. \textit{Compactness}: both actual plans are more compact than
every plan in the ensemble---more geometrically rounded, not the elongated, snaking shapes that
characterise partisan gerrymandering in the American literature.

Together, these results strongly imply that the actual boundaries were drawn with explicit
reference to URA planning area lines as an organising constraint, producing constituencies that are
administratively coherent and geographically compact. This is the opposite of what partisan
boundary manipulation looks like. The GRC system's seat-magnification effects come from the
block-vote mechanism, not from the shape of the boundaries.

\medskip
\noindent\textit{Caveat.} The GRC vs.\ SMC institutional mismatch is not a minor limitation. The
ensemble generates equal-population single-member units; the actual plans have multi-member
constituencies of different sizes. Direct metric comparisons should be interpreted cautiously, and
the results have been verified with two independent MCMC seeds (seeds 1 and 2), which produce
consistent percentile rankings across all four metrics, confirming adequate ensemble mixing.

% ============================================================
\section{The Minority Representation Rationale}

The official justification for the GRC system is ethnic minority representation. By requiring each
GRC team to include a minority candidate, the government argues, the system ensures that minority
communities are represented in parliament even in areas where they do not form a local majority.

If this rationale were the primary driver of GRC placement and sizing, we would expect GRCs to be
placed in areas with higher minority concentrations. We test this directly (Table~\ref{tab:minority}).

\begin{table}[h]
\centering
\caption{Minority population (Malay $+$ Indian) in GRC vs.\ SMC constituencies.}
\label{tab:minority}
\begin{tabular}{lccc}
\toprule
Year & GRC mean minority \% & SMC mean minority \% & $t$-test $p$-value \\
\midrule
2020 & 22.4\% & 20.5\% & 0.441 \\
2025 & 22.7\% & 18.9\% & 0.117 \\
\bottomrule
\end{tabular}
\end{table}

The difference between GRC and SMC minority percentages is not statistically significant in either
year. GRCs are not placed in areas with meaningfully higher minority populations. Additional tests
find no significant relationship between minority percentage and GRC size ($r = +0.214$, $p =
0.16$ in 2020), and no relationship between minority percentage and opposition vote share.

To be precise about what this does and does not show: the finding closes off the ethnic rationale
as the \textit{primary empirical driver} of where GRCs are placed and how large they are. It does
not prove that ethnic representation is not an outcome of the GRC system---minority candidates do
in fact win seats under GRC rules. It means that the system's structure cannot be explained by the
ethnic logic alone, because GRCs are not concentrated in minority-heavy areas.

% ============================================================
\section{Class Politics and Electoral Geography}

One of the cleanest empirical findings from the 2020 election is a strong class-politics gradient.
Using HDB flat type as a proxy for socioeconomic status---a reasonable proxy in Singapore's
context, where flat type is a direct function of income at point of purchase---we find the
following correlations with opposition vote share across the 31 contested 2020 constituencies
(Table~\ref{tab:class}).

\begin{table}[h]
\centering
\caption{Correlations with opposition vote share, 2020 contested constituencies ($n=31$).}
\label{tab:class}
\begin{tabular}{lcc}
\toprule
Variable & Pearson $r$ & $p$-value \\
\midrule
\% 4-room HDB        & $+0.483$ & 0.006 \\
\% Small HDB ($\leq$3-room) & $-0.366$ & 0.043 \\
\% Indian            & $-0.302$ & 0.099 \\
\% Minority (Malay $+$ Indian) & $-0.174$ & 0.349 \\
\bottomrule
\end{tabular}
\end{table}

Middle-class areas---constituencies with higher proportions of 4-room flats---supported the
opposition significantly more. Working-class areas supported the PAP more. Minority population had
no significant effect. This is a coherent sociological signal: the HDB flat-type gradient tracks
economic security, and voters with lower economic security may have stronger incentive to support
an incumbent party that controls public housing, welfare, and employment support.

By 2025, this signal had entirely vanished. Across all 30 contested 2025 constituencies, no
demographic variable achieves statistical significance. The highest correlation is 4-room HDB at
$r = +0.32$ ($p = 0.084$), a notable weakening from $r = +0.483$.

There are two plausible explanations for the vanishing signal, and they are not mutually
exclusive. The first is electoral: PAP's vote share rose broadly across Singapore in 2025,
possibly reflecting a rally-round-the-government response to economic uncertainty or geopolitical
tension, which would compress variance in vote shares and weaken any underlying correlation. The
second is structural: boundary changes redistributed the class concentrations that produced the
2020 signal. If middle-class HDB subzones were moved from competitive constituencies into safe or
walkover ones, the remaining contested constituencies would look demographically more homogeneous.

We cannot definitively adjudicate between these explanations with available data. What we can say
is that the class signal was real in 2020, it is gone in 2025, and the boundary changes we
document in Section~\ref{sec:permutation} are at least consistent with a structural redistribution
of class geography.

% ============================================================
\section{The Boundary Change Permutation Test}
\label{sec:permutation}

This is the headline finding.

\subsection{The question}

114 of Singapore's 332 subzones changed their constituency allocation between the 2020 and 2025
elections. A subzone is the basic unit of Singapore's urban planning geography---a named precinct
of roughly a few thousand to sixty thousand residents, drawn by the Urban Redevelopment Authority
for administrative purposes.

Of these 114 changed subzones, 52 came from constituencies where the PAP's vote share was below
55\% in 2020---constituencies where the outcome was genuinely competitive by any reasonable
standard. We call these ``competitive-origin'' subzones. In 2020, the competitive PAP
constituencies (PAP $<$55\%) were Bukit Batok SMC (54.8\%), East Coast GRC (53.4\%), and West
Coast GRC (51.7\%).

Of the same 114 changed subzones, 7 ended up in constituencies where the PAP's vote share was
below 55\% in 2025---the new competitive seats. In 2025, only one PAP constituency crossed this
threshold: Jalan Kayu SMC (51.5\%), where the Workers' Party came within 1.5 percentage points of
winning.

The question: how many of the 52 competitive-origin subzones ended up in competitive-destination
constituencies?

\subsection{The null distribution}

Under random redistribution---if the 114 changed subzones were assigned to their 2025
constituencies without reference to the political characteristics of their origins---we would
expect the overlap between competitive origins and competitive destinations to follow a
hypergeometric distribution. The expected value is:
%
\begin{equation}
  E[\text{overlap}] = \frac{52 \times 7}{114} = 3.19
\end{equation}
%
In other words, if the redistribution were random, roughly 3 of the competitive-origin subzones
would be expected to land in competitive 2025 constituencies by chance alone.

\subsection{The actual result: zero}

The actual redistricting produced zero overlap (Table~\ref{tab:contingency}).

\begin{table}[h]
\centering
\caption{Competitive-origin vs.\ competitive-destination overlap among 114 changed subzones.}
\label{tab:contingency}
\begin{tabular}{lccc}
\toprule
 & \makecell{Competitive\\2025 dest.} & \makecell{Non-competitive\\2025 dest.} & Total \\
\midrule
Competitive 2020 origin   & \textbf{0} & 52 & 52 \\
Non-competitive origin    & 7          & 55 & 62 \\
\midrule
Total                     & 7          & 107 & 114 \\
\bottomrule
\end{tabular}
\end{table}

Not one of the 52 subzones drawn from competitive 2020 constituencies ended up in a competitive
2025 constituency. Fisher's exact test (one-sided): $p = 0.012$. Hypergeometric CDF $P(X \leq 0)$:
$p = 0.012$. The actual redistricting lies at the 1.2nd percentile of the null distribution for
this statistic---a result that would occur by chance approximately 1.2\% of the time.

\begin{figure}[h]
  \centering
  \includegraphics[width=0.85\textwidth]{combined_summary.png}
  \caption{Hypergeometric null distribution with the actual result ($X=0$) marked, alongside
  seat geography panels showing competitive constituency locations in 2020 and 2025. The
  actual result falls at the 1.2nd percentile of the distribution.}
  \label{fig:combined}
\end{figure}

\subsection{Where the new competitive seat came from}

The one competitive PAP constituency in 2025 that received changed subzones is Jalan Kayu SMC.
All 7 subzones feeding into Jalan Kayu came from Ang Mo Kio GRC---which the PAP won with 71.9\%
of the vote in 2020. A PAP stronghold.

The largest single transfer was FERNVALE, a subzone of 58,800 residents---the single largest
subzone move of the entire redistricting exercise. Fernvale was extracted from Ang Mo Kio GRC and
placed into the newly formed Jalan Kayu SMC. Jalan Kayu subsequently became the closest-fought
seat in the 2025 election, with the Workers' Party winning 48.53\% of the vote.

The new competitive seat was, in effect, manufactured from stronghold population. This is not
intrinsically suspicious---new constituencies do need to come from somewhere, and somewhere must
have been a non-competitive source. But the pattern becomes meaningful when set against what
happened to the populations of the three dissolved competitive constituencies.

The competitive constituencies from 2020 were not dissolved into Jalan Kayu. They were absorbed
into safe and stronghold territory:
%
\begin{itemize}\setlength\itemsep{0.2em}
  \item \textbf{West Coast GRC} (PAP 51.7\% in 2020): largely merged into West
    Coast--Jurong West (PAP 60.0\% in 2025) and other safe GRCs.
  \item \textbf{Bukit Batok SMC} (PAP 54.8\% in 2020): dissolved; Bukit Batok Central
    subzone moved into Jurong East--Bukit Batok GRC (PAP 76.7\% in 2025).
  \item \textbf{East Coast GRC} (PAP 53.4\% in 2020): survived by name but its PAP share
    rose to 58.7\% in 2025, moving from marginal to safe.
\end{itemize}
%
The populations of the competitive seats---the communities where the election had actually been
close---went to constituencies where the PAP's margin is comfortable or overwhelming. The
populations of the PAP's safest territory seeded the one new competitive seat.

\subsection{The relocation of competitive electoral geography}

The scale of the competitive-geography shift is worth examining directly. In 2020, there were 7
constituencies where the PAP received less than 60\% of the vote---constituencies that could
plausibly become competitive with a modest swing. Of these:

\begin{itemize}\setlength\itemsep{0.2em}
  \item 6 were dissolved, merged, or substantially restructured by the 2025 redistricting.
  \item 1 survived by name (East Coast GRC) but moved from 53.4\% to 58.7\%.
\end{itemize}

In 2025, there are also 7 constituencies where the PAP received less than 60\% of the vote. Of
these:

\begin{itemize}\setlength\itemsep{0.2em}
  \item 6 are entirely new constituencies that did not exist in 2020 under those names or in
    those configurations.
  \item 1 survived (East Coast, now in a different form).
\end{itemize}

The competitive electoral landscape has been effectively relocated. The communities that were
competitive in 2020 are no longer competitive in 2025. The communities that are competitive in
2025 are different communities---Jalan Kayu in the northeast, Sembawang West, Punggol---that were
not competitive in 2020.

\begin{figure}[h]
  \centering
  \includegraphics[width=\textwidth]{choropleth_2020_2025.png}
  \caption{PAP vote share by constituency in 2020 (left) and 2025 (right). Darker shading
  indicates higher PAP vote share. The competitive seats have relocated geographically: the
  northeast communities that were marginal in 2020 are safe in 2025, and the new marginal
  seats are in different areas.}
  \label{fig:choropleth}
\end{figure}

\subsection{A correlation across all changed subzones}

Extending the analysis beyond the binary competitive/non-competitive threshold, we compute for
each of the 114 changed subzones the relationship between its origin constituency's PAP vote share
and the shift in PAP vote share between its origin and destination.

Pearson $r$ between origin PAP\% and $\Delta(\text{destination} - \text{origin})$ PAP\% is
$-0.444$ ($p < 0.001$).

More competitive origin constituencies are associated with larger upward shifts in PAP vote share
at the destination. The most competitive subzones moved to the most disproportionately safer seats.
This is consistent with what the binary test finds: competitive-origin subzones were systematically
redirected to non-competitive destinations, and the magnitude of the ``safety increase'' correlates
with how competitive the origin was.

% ============================================================
\section{Caveats and Limitations}

We have been careful to frame the findings above in specific, bounded terms. This section makes
the limits explicit.

\textit{The test conditions on the 2025 structure being fixed.} The boundary change permutation
test asks whether the allocation of the 114 changed subzones was anomalous, taking the 2025
constituency structure as given. It does not test whether the 2025 structure itself---which
constituencies exist, how many seats each has, where GRCs are placed---is anomalous. Testing the
GRC placement decision requires an ensemble model that can generate random variable-size
multi-member district configurations, which does not currently exist. The permutation test is a
test of one layer of the redistricting decision; it cannot see through to the layer above.

\textit{The result is sensitive to the competitive threshold.} We define ``competitive'' as PAP
vote share below 55\%. This is a reasonable threshold---it represents constituencies where a modest
swing could change the result---but it is a choice. With a tighter threshold (PAP $<$52\%), the
2020 competitive set shrinks and the test loses power. With a broader threshold (PAP $<$60\%), the
2025 competitive set expands and the zero-overlap result would likely change. Readers should treat
the specific $p$-value of 0.012 as a function of this threshold definition.

\textit{The 2020 competitive set is small: $n=3$ constituencies.} With only three competitive
constituencies in 2020, the entire competitive-origin set of 52 subzones comes from just three
electoral units. This makes the result fragile to misclassification at the margin: if East Coast's
2020 result (53.4\%) is better characterised as ``safe'' rather than ``competitive,'' the
competitive-origin set shrinks substantially.

\textit{Intent cannot be established.} The boundary change permutation test identifies a
statistically anomalous pattern, not a cause. The same pattern could be produced by administrative
population rebalancing---a genuine effort to equalise constituency sizes and maintain
planning-area cohesion---that happened, as a side effect, to relocate competitive seats
geographically. Singapore's population grew unevenly between 2020 and 2025; new housing
developments in Punggol and Sembawang may have required the addition of new northern
constituencies, which would explain why the new competitive seats are in those areas. We cannot
rule this out, and we do not attempt to.

What the permutation test establishes is that the observed pattern is unlikely to have arisen by
chance, and that it has a specific structure---competitive origins excluded from competitive
destinations---which is the kind of structure that would arise from deliberate competitive-seat
management. The test cannot distinguish between ``deliberate strategy'' and ``administrative
coincidence with the same statistical signature.''

\textit{MCMC robustness.} The primary ensemble uses seed 1 (9,000 post-burn-in ReCom steps).
Two additional seeds (seeds 1 and 2) were run as robustness checks; all three produce identical
percentile rankings across all four metrics (0th/100th), confirming adequate chain mixing. An
earlier run (seed 42) exhibited near-zero variance in the HDB towns-split metric, indicating a
stuck chain; results reported here use the well-mixed seed-1 ensemble.

\textit{Subzone-level analysis misses within-constituency variation.} Our demographic data is at
the subzone level; we do not have precinct-level data. This limits statistical power and means we
cannot observe within-constituency demographic heterogeneity.

% ============================================================
\section{Conclusion}

This paper has examined Singapore's 2025 redistricting from three angles: an MCMC ensemble
comparison, a demographic analysis, and a boundary change permutation test.

The MCMC results tell the clearest story about what Singapore's redistricting is \textit{not}. The
actual constituency boundaries are extreme outliers in planning-area cohesion and geometric
compactness---but they are outliers in a direction consistent with administrative logic, not
partisan boundary manipulation. Singapore does not have a gerrymander in the conventional sense of
bizarrely shaped districts optimised to concentrate or dilute votes. Its boundaries follow planning
areas. The large seat bonuses the PAP receives from its vote share come from the GRC block-vote
mechanism, not from the geographic drawing of district lines.

The demographic analysis locates a genuine empirical signal that has since disappeared. In 2020,
middle-class HDB areas significantly predicted opposition vote share; by 2025 that signal had
vanished. Whether this reflects a uniform PAP swing across all socioeconomic groups, or a
redistribution of class geography through boundary changes, cannot be resolved with available data.

The boundary change permutation test is the new finding. Among 114 subzones that changed
constituency, those originating in competitive 2020 constituencies were systematically excluded
from competitive 2025 constituencies. Not one of 52 competitive-origin subzones ended up in a
competitive 2025 constituency, against an expected value of 3.2 under random redistribution.
Fisher's exact $p = 0.012$.

The structure of the anomaly is specific. The one new competitive seat in 2025 (Jalan Kayu) was
seeded entirely from Ang Mo Kio stronghold population. The populations of the three dissolved
competitive constituencies were absorbed into safe or walkover districts. The competitive electoral
geography was relocated: in 2025, the competitive seats are in entirely different communities from
2020.

Whether this reflects deliberate strategy or administrative coincidence, we cannot say. The data
demonstrates the pattern is statistically anomalous. It does not demonstrate intent. Both
conclusions deserve to be stated plainly.

What the findings do support is a more focused critique than ``gerrymandering.'' The primary
mechanism producing large PAP seat majorities from plurality vote shares is the GRC block-vote
system, not boundary drawing. The secondary mechanism---the one identified by the permutation
test---is the management of which communities become competitive, through the selective allocation
of changed subzones. These are separable claims. The first is well-established in the political
science literature; the second is what this paper adds.

% ============================================================
\section*{Methodological Note}
\addcontentsline{toc}{section}{Methodological Note}

\textbf{Data.} Electoral results from the Elections Department Singapore (ELD) via data.gov.sg.
Subzone geographic boundaries from the Urban Redevelopment Authority (URA) Master Plan 2019.
Census demographics from Singapore Census of Population 2020, at subzone level. HDB flat-type
counts from HDB property table. All data is publicly available.

\textbf{MCMC ensemble.} Implemented using GerryChain (Metric Geometry and Gerrymandering Group,
MIT). 9,000 post-burn-in ReCom steps, $k=33$ districts, population tolerance 10\%, primary seed 1.
Robustness verified with seed 2 (9,000 steps); both seeds produce identical percentile rankings.
Subzone adjacency graph: 332 nodes, 850 edges (rook contiguity).

\textbf{Permutation test.} Hypergeometric test and Fisher's exact test. Competitive threshold: PAP
vote share $<$55\% in the relevant election year. All computations in Python using
\texttt{scipy.stats}.

\textbf{Code and data.} Available at \url{https://github.com/davidcagoh/sg-redistricting-2025}.

% ============================================================
\begin{thebibliography}{99}

\bibitem[Herschlag et~al.(2020)]{herschlag2020}
Herschlag, G., Kang, H.~S., Luo, J., Graves, C.~V., Bangia, S., Ravier, R.,
  \& Mattingly, J.~C. (2020).
\newblock Quantifying gerrymandering in North Carolina.
\newblock \textit{Statistics and Public Policy}, 7(1), 30--38.

\bibitem[McGhee(2014)]{mcghee2014}
McGhee, E. (2014).
\newblock Measuring partisan bias in single-member district electoral systems.
\newblock \textit{Legislative Studies Quarterly}, 39(1), 55--85.

\bibitem[Singapore Elections Department(2020, 2025)]{eld2025}
Singapore Elections Department. (2020, 2025).
\newblock \textit{Parliamentary General Election results}.
\newblock data.gov.sg.

\bibitem[Urban Redevelopment Authority(2019)]{ura2019}
Urban Redevelopment Authority. (2019).
\newblock \textit{Master Plan 2019 Subzone Boundary}.
\newblock data.gov.sg.

\end{thebibliography}

\end{document}
